\scanpage{102}
Les \deleted{objets} faisceaux de cohomologie de $L_{\bullet}\overset{L}{\otimes}M_{\bullet}$ sont appelés \uline{hypertor} de $L_{\bullet},M_{\bullet}$, notation $\mathbb T\mathrm{or}_i^{\underline O}(L_{\bullet},M_{\bullet})$. \deleted{\unsure{Si $L_{\bullet}$ ou $M$|Si $L_{\bullet},M$}} Comme il s'agit ici de foncteurs \inserted{``exacts''} de catégories triangulées, tout triangle exact en l'un des arguments donne lieu à la suite exacte infinie habituelle. En particulier, les $\mathbb T\mathrm{or}_i(\ ,M_{\bullet})$ définissent un foncteur cohomologique en le premier argument, lorsqu'on prend pour celui-ci un Module ordinaire.

La formation du $\deleted{\unsure{K}}L_{\bullet}\overset{\mathbb L}{\otimes}M_{\bullet}$ commute au changement de base \uline{plat} \inserted{$f:X\longrightarrow Y$}, quand celui-ci est défini \unsure{naïvement} : \deleted{\unsure{[mais alors}}
\[
\boxed{\mathbb L f^*(M_{\bullet}\overset{\mathbb L}{\otimes}L_{\bullet})\simeq\mathbb L f^*(M_{\bullet})\overset{\mathbb L}{\otimes}\mathbb L f^*(L_{\bullet})}
\]
Pour définir le foncteur \unsure{$f\mathbin{-}$|$f^*$ ou} $\mathbb L f^*$ entre $D^-$, il faut noter que la \uline{platitude} de $f$ implique que $f^*$ transforme quasi-isomorphismes en quasi-isomorphismes. D'autre part, $f^*$ transforme plats en plats, d'où la commutation.
% slanted
\margin{est-ce vrai dans le contexte des topos généraux ?}
--- En particulier, la formation des $L_{\bullet}\overset{\mathbb L}{\otimes}M_{\bullet}$ commute au passage aux fibres, ce qui dans le cas qui nous occupe permettra \inserted{\unsure{pfois|parfois|tjours}} un calcul pratiquement trivial des $\overset{\mathbb L}{\otimes}$.
% archive sheet number: 101
